\section{Introduction}
\label{sec:ch5-introduction}
Interactive systems are increasingly embedded in indoor environments, leveraging sensing, actuation, and computation to regulate environmental parameters like lighting, temperature, and sound~\cite{alavi2019introduction, alavi2017comfort}. These systems enhance comfort, efficiency, and adaptability by reacting to changing conditions~\cite{nguyen2022towards, nguyen2024adaptive, deng2019diffusive}. This reflects an optimisation-oriented paradigm where environmental qualities are variables to be measured, predicted, or personalised~\cite{heidegger2006building}.

However, this perspective neglects deeper experiential dimensions rooted in our innate connection to nature~\cite{wilson1986biophilia, kellert1995biophilia, kellert2018nature}. Biophilic design, grounded in this affinity, embeds sensory, material, and ecological references into the built environment~\cite{wilson1986biophilia, kellert1995biophilia}. Such experiences are characterised by ecological connectedness, multisensory engagement beyond vision, and environmental variability~\cite{rao2025designing}, and are linked to stress reduction, attentional restoration, and emotional regulation~\cite{frumkin2017nature, hartig2014nature, lefosse2023biophilia}. Yet these dimensions remain under-represented in current paradigms~\cite{rao2026emotion}, which prioritise stability over sensory and processual engagement. This creates a fundamental tension: optimisation reduces variability, whereas these experiential qualities depend on it.

Rather than maintaining predefined optimal states, biophilic design foregrounds ongoing sensory engagement and environmental variation, positioning variability and ecological processes as central to interaction~\cite{rao2025designing, rao2025lessons}.

These ideas are deeply embedded in architectural traditions. For example, Haveli architecture in South Asia organises homes around a central courtyard that brings light, air, and weather into the interior. Indoor experience is therefore shaped through variation, requiring occupants to continually adjust movement, position, and duration of stay in response to changing conditions~\cite{verma2022appraisal}. This illustrates how biophilic design operates through spatial organisation, material choices, and environmental variation to shape how occupants sense and act. Such qualities resist reduction to discrete parameters and remain difficult to articulate within optimisation-oriented models. This gap between experience-oriented architectural practice and its limited articulation in interactive systems is the focus of this paper. While biophilic design also encompasses more-than-human considerations, we focus here on human sensory and embodied interaction in indoor environments.

Despite this architectural grounding, interactive systems have only partially explored how biophilic qualities might inform interaction within human–computer interaction (HCI). This agenda is articulated through human--building interaction (HBI), which examines how building technologies shape everyday experience at architectural scale~\cite{alavi2019introduction, becerik2022field}. HBI extends experience-centred perspectives to technology-mediated environments, attending to occupants, spatial systems, and environmental mediation. While HBI has broadened inquiry beyond comfort toward aesthetics~\cite{economidou2021moving}, social experience~\cite{hoare2015hide}, and place-making~\cite{beale2022digital}, biophilic qualities remain largely implicit or isolated. At the same time, biophilic and experiential perspectives are gaining momentum across HCI and HBI~\cite{rao2025lessons, rao2025designing, margariti2024human, margariti2024evaluating, han2024superarchitectural}. However, insights are often expressed through project-specific examples~\cite{behrens2018sentiment}, metaphors~\cite{deng2019diffusive}, or singular concepts~\cite{wang2022nature}, limiting cumulative knowledge-building.

Through expert workshops with practising architects and designers, we elicited how biophilic experiences are imagined, reasoned about, and translated into interactive spatial propositions. Engaging experts whose work centres on materiality, atmosphere, and human–environment relations enabled access to experiential reasoning (e.g. temporality, spatial behaviour, and atmospheric composition) that are difficult to surface through technology-led design processes alone~\cite{alavi2019introduction, alavi2018artifacts}. 

Rather than producing finished designs, this process aimed to surface and structure this reasoning into a \textbf{biophilic interaction design framework for indoor environments.} We ask:

\begin{itemize}
    \item \textbf{RQ1}: What biophilic experiences and interactions are being imagined for future indoor environments?
    \item \textbf{RQ2}: Why do these designs matter, and what values or experiential orientations underpin them?
    \item \textbf{RQ3}: How can these experiences inform interaction forms for bodily interaction with interactive indoor environments?
\end{itemize}

We address these questions through three speculative design sessions (S1---S3): Session~1 surfaces missing experiences, Session~2 explores spatial translation, and Session~3 focuses on sound as a modality. This design-led approach surfaces expert imaginaries and reasoning often overlooked in technology-led discourse.

This paper contributes a \textbf{biophilic interaction design framework} comprising three analytic layers: (1) dimensions that map the design space, (2) design values that articulate experiential reasoning behind expert decisions, and (3) interaction forms that describe recurring logics of bodily-spatial engagement. Together, these layers provide an integrated framework for analysing and designing biophilic interactions in indoor environments. \textbf{To our knowledge, this is the first structured design framework for biophilic interaction in HCI.} Our framework supports an expanded vision of interaction design shaped by bodily inhabitation, environmental mediation, and situated sensory experience.